%%%%%ENTREES%%%%
\begin{tabarg}
\multicolumn{3}{l}{\titretab {Entr\'ees}}\\
\hline
$X_{in}$ & Vecteur, \texttt{float} & Variable d'origine\\
$time_{in}$ & Vecteur, \texttt{double} & Base de temps de $X_{in}$ (secondes)\\
$toler\_mnq$ & Coeff., \texttt{Int\_32} & Seuil de calcul de la mvariance (nb. �chant). A chaque instant de la base de temps d'arriv�e, 
la variance est calcul�e si le nombre d'�chantillons dans l'intervalle de temps est strictement sup�rieur � $toler\_mnq$. Sinon, 
le r�sultat est "valeur manquante". Une valeur de 0 provoque le calcul dans tous les cas (m�me s'il n'y a qu'un seul �chantillon).\\
$time_{out}$ & Vecteur, \texttt{double} & Base de temps d'arriv�e (secondes)\\
$time\_bnds$ & Coeff, \texttt{string} & cha�ne d�crivant les bornes de temps selon la norme AIDA (dans fvb, utiliser \$Arg4.time\_bnds)\\
$size_{out}$ & Coeff., \texttt{Int\_32} & Nombre de valeurs dans la base de temps 
d'arriv�e (dans fvb, utiliser \$Arg4.NbVal)\\
\end{tabarg}

\medskip
%%%%%SORTIES%%%%
\begin{tabarg}
\multicolumn{3}{l}{\titretab {Sorties}}\\
\hline
$V_{x}$ & Vecteur, \texttt{float} & Variance de $X_{in}$ �chantillonn�e sur
la base de $time_{out}$ (m�me unit� que $X_{in})$\\
\end{tabarg}
